\chapter{Global Types as Grammars}
\label{ch:grammar}

\index{grammar!multiparty}
\index{grammar!context-free}
\index{global type}
A global session type is a grammar. Its terminal symbols are messages, each with
a sender and a receiver; its productions are the type's sequence, choice,
parallel composition and recursion; the language it generates is the set of
legal conversations. LeanFM writes global types this way so that
requirements can be stacked: every production can be read as if it were the
top-level protocol, and a larger protocol is built by referring to smaller
ones. \textbf{Proposed notation:} the \texttt{Grammar}
algebra and BNF below are not an executable API. Durable artifacts use
\texttt{GrammarExpr} (
\texttt{event}, \texttt{seq}, \texttt{choice}, \texttt{parallel},
\texttt{thresholdJoin}, \texttt{guard}, \texttt{ref}, \texttt{repeat}).
Current validation checks grammar/task and message references, atom endpoints,
and threshold bounds; it does not compile this notation into trace generation or
recognition, or prove actor projection. The queue rules later in this chapter
are model semantics, not a proof of equivalence to a Go runtime.

\section{A small authoring algebra}

\begin{lstlisting}
inductive Grammar (Event : Type) where
  | empty
  | event  : Event -> Grammar Event
  | seq    : Grammar Event -> Grammar Event -> Grammar Event
  | choice : List (Grammar Event) -> Grammar Event
  | par    : Grammar Event -> Grammar Event -> Grammar Event
  | where_ : (List Event -> Bool) -> Grammar Event -> Grammar Event
\end{lstlisting}

Each constructor is a session-type operator. \texttt{event} is one message
from a sender to a receiver. \texttt{seq} preserves time order. \texttt{choice} exposes alternatives rather
than burying them in control flow. \texttt{par} admits interleavings while
preserving the internal order of each side. A \texttt{where\_} clause checks a
fact accumulated from the prior trace and provides the escape from context-free
structure when genuine context sensitivity is needed.

\index{choice!directed}
A choice does not need a separate annotation naming who makes it, as Scribble's
\texttt{choice at} does. Each branch
eventually resolves to terminal events, and each terminal resolves to framed
bytes with a source and destination. The first observed terminal that
distinguishes two branches identifies both the selected branch and the actor
that emitted the distinguishing bytes. If two branches resolve to the same
observable bytes, then at this abstraction they are not distinct choices and
should be merged. This removes an otherwise ambiguous second account of agency
from the grammar.

One condition makes this sound: the distinguishing first messages of all
branches must be sent by the \emph{same} actor. If Alice's message starts one
branch and Bob's starts the other, nobody can implement the choice locally;
both may act at once and the system leaves the grammar. This is the
non-local-choice (realizability) problem of multiparty protocols. In
session-type terms the choice must be \emph{directed}: an internal choice of
one role, observed by the others as an external choice on their queues. A
second condition goes with it: every role whose later behaviour differs between
branches must receive a message that tells it which branch was taken, or must
have a declared timer that ends its part of the session
(Chapter~\ref{ch:vision}). LeanFM does not yet perform either validation.

Simple finite conversations can compile to regular languages. Nested request
structure motivates context-free productions. Correlation, freshness, uniqueness,
and authorization typically require parameter flow or guarded productions.

\section{A BNF for composable interactions}

The authoring notation can be presented independently of Lean as a small BNF.
Named productions are interactions; referring to a production composes that
interaction into a larger one.

\begin{lstlisting}[language={}]
<specification> ::= <production>
                  | <production> <specification>

<production>   ::= <name> "::=" <expression>

<expression>   ::= "empty"
                  | <event>
                  | <name>
                  | "(" <expression> ";" <expression> ")"
                  | "(" <expression> "|" <expression> ")"
                  | "(" <expression> "||" <expression> ")"
                  | "repeat" "(" <expression> ")"
                  | "join" "(" <natural> "of" <expressions> ")"
                  | "when" "(" <predicate> "," <expression> ")"
                  | "after" "(" <name> "," <expression> ")"

<event>        ::= <actor> "->" <actor> ":"
                   <message> "(" <fields> ")"

<fields>       ::= ""
                  | <field>
                  | <field> "," <fields>

<field>        ::= <name> "=" <value>
\end{lstlisting}

The operators have the same meanings as the grammar AST: semicolon is sequence,
vertical bar is alternative, and double vertical bar is parallel composition by
order-preserving interleaving. \texttt{after(P,E)} includes \texttt{E} only on a
branch in which production \texttt{P} completed. Predicates and repeated field
variables supply the context-sensitive correlation that plain context-free syntax
does not express by itself.

A threshold join is explicit rather than inferred from the number of observed
predecessors:

\begin{lstlisting}[language={}]
firstTwoReplicas ::= join(2 of replicaA, replicaB, replicaC)
firstTwoThenReply ::= (firstTwoReplicas ; reply)
\end{lstlisting}

The continuation becomes enabled after any two of the three branches complete.
Its first emitted event records \texttt{joinRequired=2},
\texttt{joinTotal=3}, and exactly two selected branch-terminal IDs. Those IDs
also occur in the event's ordinary \texttt{prior} list, which may contain other
backpointers such as the task start. The threshold declaration says why two
predecessors suffice; the selected IDs say which two satisfied this instance.
Unselected branches are not implicitly canceled and may still emit events.
Cancellation, if required, must be an explicit message or production.

The OAuth example decomposes into independently readable interactions.  These
generated views show the two complete branches without making a message arrow
look like logical implication:

\begin{figure}[htbp]
\centering
\AccessibleFigure{OAuth granted interaction. Message from Browser to Application: Application Request. Message from Application to Browser: Authorization Redirect. Message from Browser to Authorization: Authorization Request. Message from Authorization to Browser: Authorization Code Returned. Message from Browser to Application: Authorization Callback. Message from Application to Authorization: Token Request. Message from Authorization to Application: Access Token Issued. Message from Application to Resource: Resource Request. Message from Resource to Application: Resource Returned. Message from Application to Browser: Application Response.}{%
\resizebox{\textwidth}{!}{% Generated by LeanFM.MDPDiagram.renderInteractionTeX.
\begin{tikzpicture}[actor/.style={draw=diagramink,text=diagramink,fill=diagramsurface,rounded corners,minimum width=2.3cm,minimum height=7mm,font=\small},message/.style={->,>=Stealth,draw=leanblue,thick}]
\node[actor] (a0) at (0,0) {Browser};
\draw[dashed,draw=diagramink] (a0.south) -- ++(0,-11);
\node[actor] (a1) at (6,0) {Application};
\draw[dashed,draw=diagramink] (a1.south) -- ++(0,-11);
\node[actor] (a2) at (12,0) {Authorization};
\draw[dashed,draw=diagramink] (a2.south) -- ++(0,-11);
\node[actor] (a3) at (18,0) {Resource};
\draw[dashed,draw=diagramink] (a3.south) -- ++(0,-11);
\draw[message] (0,-1) -- node[above,fill=pagebackground,text=diagramink,font=\scriptsize] {ApplicationRequest(resource)} (6,-1);
\draw[message] (6,-2) -- node[above,fill=pagebackground,text=diagramink,font=\scriptsize] {AuthorizationRedirect(state; callback)} (0,-2);
\draw[message] (0,-3) -- node[above,fill=pagebackground,text=diagramink,font=\scriptsize] {AuthorizationRequest(state; callback)} (12,-3);
\draw[message] (12,-4) -- node[above,fill=pagebackground,text=diagramink,font=\scriptsize] {AuthorizationCodeReturned(state; code)} (0,-4);
\draw[message] (0,-5) -- node[above,fill=pagebackground,text=diagramink,font=\scriptsize] {AuthorizationCallback(state; code)} (6,-5);
\draw[message] (6,-6) -- node[above,fill=pagebackground,text=diagramink,font=\scriptsize] {TokenRequest(code)} (12,-6);
\draw[message] (12,-7) -- node[above,fill=pagebackground,text=diagramink,font=\scriptsize] {AccessTokenIssued(token)} (6,-7);
\draw[message] (6,-8) -- node[above,fill=pagebackground,text=diagramink,font=\scriptsize] {ResourceRequest(token)} (18,-8);
\draw[message] (18,-9) -- node[above,fill=pagebackground,text=diagramink,font=\scriptsize] {ResourceReturned(profile)} (6,-9);
\draw[message] (6,-10) -- node[above,fill=pagebackground,text=diagramink,font=\scriptsize] {ApplicationResponse(profile)} (0,-10);
\end{tikzpicture}
}}
\caption{Granted OAuth interaction. The resource exchange occurs only after the
authorization-code branch succeeds.}
\label{fig:oauth-grammar-granted}
\end{figure}

\begin{figure}[htbp]
\centering
\AccessibleFigure{OAuth denied interaction. Message from Browser to Application: Application Request. Message from Application to Browser: Authorization Redirect. Message from Browser to Authorization: Authorization Request. Message from Authorization to Browser: Authorization Error Returned. Message from Browser to Application: Authorization Error Callback.}{%
\resizebox{\textwidth}{!}{% Generated by LeanFM.MDPDiagram.renderInteractionTeX.
\begin{tikzpicture}[actor/.style={draw=diagramink,text=diagramink,fill=diagramsurface,rounded corners,minimum width=2.3cm,minimum height=7mm,font=\small},message/.style={->,>=Stealth,draw=leanblue,thick}]
\node[actor] (a0) at (0,0) {Browser};
\draw[dashed,draw=diagramink] (a0.south) -- ++(0,-6);
\node[actor] (a1) at (6,0) {Application};
\draw[dashed,draw=diagramink] (a1.south) -- ++(0,-6);
\node[actor] (a2) at (12,0) {Authorization};
\draw[dashed,draw=diagramink] (a2.south) -- ++(0,-6);
\node[actor] (a3) at (18,0) {Resource};
\draw[dashed,draw=diagramink] (a3.south) -- ++(0,-6);
\draw[message] (0,-1) -- node[above,fill=pagebackground,text=diagramink,font=\scriptsize] {ApplicationRequest(resource)} (6,-1);
\draw[message] (6,-2) -- node[above,fill=pagebackground,text=diagramink,font=\scriptsize] {AuthorizationRedirect(state; callback)} (0,-2);
\draw[message] (0,-3) -- node[above,fill=pagebackground,text=diagramink,font=\scriptsize] {AuthorizationRequest(state; callback)} (12,-3);
\draw[message] (12,-4) -- node[above,fill=pagebackground,text=diagramink,font=\scriptsize] {AuthorizationErrorReturned(state; denied)} (0,-4);
\draw[message] (0,-5) -- node[above,fill=pagebackground,text=diagramink,font=\scriptsize] {AuthorizationErrorCallback(state; denied)} (6,-5);
\end{tikzpicture}
}}
\caption{Denied OAuth interaction. No token or resource request follows the error.}
\label{fig:oauth-grammar-denied}
\end{figure}
\clearpage

Chapter~\ref{ch:oauth-example} gives the complete production definitions,
nested alternatives and guarded exchange; these terminal-level diagrams isolate
the granted and denied outcomes for comparison.

\section{Parameters carry information}

\index{parameter flow}
A production should bind values at one event and reuse them later. Conceptually:

\begin{lstlisting}
def login (browser server provider : Actor) (state : OAuthState) : Grammar Event :=
  event (redirect server browser provider state) >>>
  event (authorize browser provider state) >>>
  callback browser server state
\end{lstlisting}

The repeated \texttt{state} is not decorative. It states the information-flow
relationship that distinguishes a matching callback from an unrelated one. Lean's
types, constructors, and ordinary function parameters are the surface language;
the fixed library supplies grammar composition and semantics.

\section{Projection: from the global type to each role}

\index{projection}
\index{local type}
A participant's implementation need not know the global execution state. The
global type states the contract, and projection derives each role's
\emph{local type}: the same grammar with every message the role neither sends
nor receives erased, and every remaining message rewritten as a send
(\texttt{!}) or a receive (\texttt{?}). For the granted OAuth branch of
\cref{fig:oauth-grammar-granted}, the resource server's local type is two lines
long:

\begin{lstlisting}[language={}]
Resource ::= ? Application : resourceRequest(token) ;
             ! Application : resourceReturned(resource)
\end{lstlisting}

The resource server never learns how the token was obtained, and its
implementation should not depend on it. Projection is defined only for a
well-formed global type: where a choice is erased, the branches a role still
sees must be the same, or must start with different messages that role
receives. Composition should preserve independently specified modules and make
shared messages explicit. Proving that projection and composition preserve the
global type's properties is a core research obligation, not an implementation
footnote; LeanFM does not yet compute projections (a \textbf{future
interface}).

\section{Declared physics}

\index{physics of a channel}
\index{channel kind}
A global type says which message may follow which. It does not say what
happens to a message between its send and its receive. LeanFM requires that to
be declared beside the protocol, as the \emph{physics} of each channel
(\textbf{proposed notation}):

\begin{lstlisting}[language={}]
physics Transfer where
  channel Sender -> Receiver : udp (capacity := 2, duplication := never)
  channel Receiver -> Sender : udp (capacity := 2, duplication := never)
  timer Sender : mayFireEarly
\end{lstlisting}

Like the grammar, the declaration is ordinary Lean data. The six axes of
Chapter~\ref{ch:vision} are six fields, and a preset is a value of the
structure that a declaration updates:

\begin{lstlisting}
inductive WhenFull  | block | reject | dropNewest | dropOldest
inductive Ordering  | fifo | none
inductive Loss      | never | mayDrop
inductive Duplication | never | mayDuplicate

structure Sharding (Msg : Type) where
  key        : Msg -> String     -- a field of the message
  partitions : Nat

structure Channel (Msg : Type) where
  capacity    : Nat
  whenFull    : WhenFull
  ordering    : Ordering         -- within one partition
  sharding    : Option (Sharding Msg)
  loss        : Loss
  duplication : Duplication
  capacityPos : 0 < capacity

def udp (n : Nat) (h : 0 < n) : Channel Msg :=
  { capacity := n, capacityPos := h, whenFull := .dropNewest,
    ordering := .none, sharding := none,
    loss := .mayDrop, duplication := .mayDuplicate }
\end{lstlisting}

Per-key ordering is not a separate value: it is \texttt{fifo} ordering on a
sharded channel, where each partition is its own FIFO queue.

The checked system is then the product of the roles' local types and their
channels. Each field contributes, or removes, transitions of the product:

\begin{itemize}
  \item \texttt{capacity} and \texttt{whenFull}: a send to a full queue blocks
        the sender, delivers a refusal message back to it, or discards a
        message;
  \item \texttt{ordering}: under \texttt{fifo} only the head of a queue may be
        delivered; under \texttt{none} any queued message may be;
  \item \texttt{sharding}: the channel is \texttt{partitions} queues, each with
        the declared capacity, and a send goes to the queue its key selects;
        for checking, keys are abstracted to as many distinct values as the
        protocol's messages can tell apart;
  \item \texttt{loss}: a transition that removes any queued message;
  \item \texttt{duplication}: a transition that delivers a queued message and
        leaves it queued.
\end{itemize}

A timer contributes a transition that may fire whenever its role is waiting.
Every one of these is a choice the checker explores, so each \texttt{may}
multiplies the state space; the presets keep a declaration honest without
making it long.
Chapter~\ref{ch:tftp-physics} builds this product for TFTP and shows a global
type that is well formed in the classic sense, and whose original
implementation duplicates every block once a single timer fires early.

The rest of this chapter defines the default physics, the one that LeanFM's
actor runtime assumes when nothing more is declared: bounded, reliable,
blocking queues.

\section{Local types run on actors with bounded queues}

\index{session identifier}
Each role is played by an actor, and one actor may take part in many sessions
at once. An actor has separate positive finite-capacity inbound and outbound message queues. Queue
capacity is part of the observable execution model: producing a message may block
when the outbound queue is full, and transport may block when the destination's
inbound queue is full. Sending is therefore distinct from delivery and from the
destination actor consuming the message.

Messages carry a session identifier, naming the instance of the global type
they belong to, as well as a task identifier. Their pair is
the correlation key for actor-local work:

\begin{lstlisting}
structure WorkKey where
  session : SessionId
  task    : TaskId

structure ActorRuntime where
  inboundCapacity  : Nat
  outboundCapacity : Nat
  inbound          : List Envelope
  outbound         : List Envelope
  inFlight         : List (Prod WorkKey GrammarResidual)
  inboundBounded   : inbound.length <= inboundCapacity
  outboundBounded  : outbound.length <= outboundCapacity
\end{lstlisting}

An actor may have zero or many distinct \texttt{WorkKey}s in flight. Consuming an
inbound envelope selects or creates the residual associated with that envelope's
\texttt{(session, task)} pair; it does not replace one actor-wide ``current task.''
Messages from different keys may interleave while the grammar preserves the order
required within each conversation.

The queues have Go-channel-like operational semantics. Capacity is enforced when
an operation is attempted, not checked only after constructing a new state:

\begin{lstlisting}
send(channel, message) =
  if channel.length < channel.capacity
  then sent(channel ++ [message])
  else blocked

receive(channel) =
  if channel = message :: rest
  then received(message, rest)
  else blocked

tryReceive(channel) =
  if channel = message :: rest
  then some(message, rest)
  else none
\end{lstlisting}

A blocked operation makes no progress and leaves the channel unchanged. A
blocking actor may run again when the required space or data becomes available.
\texttt{tryReceive} is a different operation: \texttt{none} reports that the
channel is empty now and lets the actor continue. It does not put the actor to
sleep, and it must not be confused with receiving a valid optional payload.

The resulting transitions are explicit:

\begin{itemize}
  \item an actor step appends a produced envelope to its outbound queue;
  \item a transport step moves the outbound head to the destination's inbound
        queue when both operations are within capacity;
  \item a receive step removes an inbound envelope and advances the matching
        \texttt{(session, task)} residual;
  \item sending to a full outbound queue blocks the sender without enqueueing;
  \item transport to a full inbound queue blocks without removing the outbound
        head;
  \item blocking receive on an empty inbound queue blocks, while nonblocking
        \texttt{tryReceive} returns \texttt{none} and permits later work.
\end{itemize}

Queue order and predecessor order are not identical. FIFO determines which queued
envelope is available next, while \texttt{priorIds} records the event-order graph,
including forks and joins across actors. Incomparable branches may commute; the
pointer relation does not itself claim that one event caused another.
